\documentclass[10pt,a4paper]{amsart}
\usepackage{amsmath,amssymb,amsthm,graphicx}
\usepackage{verbatim}
\newenvironment{reference}{\comment}{\endcomment}
%\newenvironment{reference}{}{}
\newenvironment{slogan}{\comment}{\endcomment}
\newenvironment{history}{\comment}{\endcomment}
% For commutative diagrams we use Xy-pic
\usepackage[all]{xy}

% We use 2cell for 2-commutative diagrams.
\xyoption{2cell}
% Theorem environments.
%
\theoremstyle{plain}
\newtheorem{theorem}[subsection]{Theorem}
\newtheorem{proposition}[subsection]{Proposition}
\newtheorem{lemma}[subsection]{Lemma}

\theoremstyle{definition}
\newtheorem{definition}[subsection]{Definition}
\newtheorem{example}[subsection]{Example}
\newtheorem{exercise}[subsection]{Exercise}
\newtheorem{situation}[subsection]{Situation}

\theoremstyle{remark}
\newtheorem{remark}[subsection]{Remark}
\newtheorem{remarks}[subsection]{Remarks}

\numberwithin{equation}{subsection}

% Macros
%
\def\lim{\mathop{\mathrm{lim}}\nolimits}
\def\colim{\mathop{\mathrm{colim}}\nolimits}
\def\Spec{\mathop{\mathrm{Spec}}}
\def\Hom{\mathop{\mathrm{Hom}}\nolimits}
\def\Ext{\mathop{\mathrm{Ext}}\nolimits}
\def\SheafHom{\mathop{\mathcal{H}\!\mathit{om}}\nolimits}
\def\SheafExt{\mathop{\mathcal{E}\!\mathit{xt}}\nolimits}
\def\Sch{\mathit{Sch}}
\def\Mor{\mathop{\mathrm{Mor}}\nolimits}
\def\Ob{\mathop{\mathrm{Ob}}\nolimits}
\def\Sh{\mathop{\mathit{Sh}}\nolimits}
\def\NL{\mathop{N\!L}\nolimits}
\def\CH{\mathop{\mathrm{CH}}\nolimits}
\def\proetale{{pro\text{-}\acute{e}tale}}
\def\etale{{\acute{e}tale}}
\def\QCoh{\mathit{QCoh}}
\def\Ker{\mathop{\mathrm{Ker}}}
\def\Im{\mathop{\mathrm{Im}}}
\def\Coker{\mathop{\mathrm{Coker}}}
\def\Coim{\mathop{\mathrm{Coim}}}

% Boxtimes
%
\DeclareMathSymbol{\boxtimes}{\mathbin}{AMSa}{"02}

%
% Macros for moduli stacks/spaces
%
\def\QCohstack{\mathcal{QC}\!\mathit{oh}}
\def\Cohstack{\mathcal{C}\!\mathit{oh}}
\def\Spacesstack{\mathcal{S}\!\mathit{paces}}
\def\Quotfunctor{\mathrm{Quot}}
\def\Hilbfunctor{\mathrm{Hilb}}
\def\Curvesstack{\mathcal{C}\!\mathit{urves}}
\def\Polarizedstack{\mathcal{P}\!\mathit{olarized}}
\def\Complexesstack{\mathcal{C}\!\mathit{omplexes}}
% \Pic is the operator that assigns to X its picard group, usage \Pic(X)
% \Picardstack_{X/B} denotes the Picard stack of X over B
% \Picardfunctor_{X/B} denotes the Picard functor of X over B
\def\Pic{\mathop{\mathrm{Pic}}\nolimits}
\def\Picardstack{\mathcal{P}\!\mathit{ic}}
\def\Picardfunctor{\mathrm{Pic}}
\def\Deformationcategory{\mathcal{D}\!\mathit{ef}}
\usepackage{cleveref}
\crefname{originallemma}{original lemma}{original lemmas}
\Crefname{originallemma}{Original lemma}{Original lemmas}
\makeatletter
\newlabel{lemma-RHom-ext}{{lemma-RHom-ext}{source}{}{sourcepointer-1}{}}
\newlabel{lemma-RHom-ext@cref}{{[originallemma][0][]lemma-RHom-ext}{[1][0][]source}{}{}{}}
\makeatother
\crefname{originalreference}{original reference}{original references}
\Crefname{originalreference}{Original reference}{Original references}
\makeatletter
\newlabel{algebra-lemma-ext-noetherian}{{algebra-lemma-ext-noetherian}{source}{}{sourcepointer-2}{}}
\newlabel{algebra-lemma-ext-noetherian@cref}{{[originalreference][0][]algebra-lemma-ext-noetherian}{[1][0][]source}{}{}{}}
\makeatother
\crefname{originalreference}{original reference}{original references}
\Crefname{originalreference}{Original reference}{Original references}
\makeatletter
\newlabel{derived-section-triangulated-sub}{{derived-section-triangulated-sub}{source}{}{sourcepointer-3}{}}
\newlabel{derived-section-triangulated-sub@cref}{{[originalreference][0][]derived-section-triangulated-sub}{[1][0][]source}{}{}{}}
\makeatother
\title{The derived internal hom along a finite ring map preserves bounded below coherent complexes}
\author{The Stacks Project Authors}
\date{stacks-project:dualizing.tex:local-snapshot-20261012; GNU FDL 1.2 or later}
\begin{document}
\maketitle
\noindent\textit{Editorial note.} All delivered mathematical text is the author's own, copied from the pinned Stacks Project chapter source sources/stacks-dualizing/pinned-original-source.tex; no statement or proof step was rewritten, reordered or removed.\par\medskip
\noindent\textit{Editorial note.} Preamble provenance: the Stacks Project keeps one single shared preamble for every chapter; this card takes it directly from the staged sources/stacks-dualizing/preamble.tex as three verbatim spans (lines 8-12, 14-18 and 154-217), so no derived file was written and nothing was added or redefined. The wrapper class and the base packages are supplied by the assembler.\par\medskip
\noindent\textit{Editorial note.} Bibliography provenance: neither the counted statement nor its proof contains a citation (this lane rescans the whole span before writing the card), so the card carries no bibliography and no entry was reconstructed.\par\medskip
\noindent\textit{Editorial note.} Selection provenance: this card was opened under the file-based control contract (CONTROL.md), with the mechanical pre-card gate run on its statement and proof interval before the card was written.\par\medskip
\noindent\textit{Allowed standard prerequisites (editorial).}\begin{itemize}
\item The identification of the derived internal hom along a ring map with the Ext groups of the base ring. Invoked inside the counted proof (source line 2489). Turns the finiteness of the cohomology of the internal hom into the finiteness of Ext modules of the base ring.
\item The finiteness of Ext over a Noetherian ring between a finite module and a finite module. Invoked inside the counted proof (source line 2498). Supplies the finiteness of the spectral sequence's E2 terms.
\end{itemize}\medskip
\ifdefined\hypertarget\hypertarget{sourcepointer-1}{}\fi
\noindent\textit{Source reference.} Original lemma lemma-RHom-ext: sources/stacks-dualizing/pinned-original-source.tex, source lines 2452--2452. Navigation pointer only: the counted text refers to the Stacks item lemma-RHom-ext, which is an existing labeled result of this same chapter and is not part of this card. The pointer carries the target's own label as its number and its definition line as its span; its content is not reproduced here.\par\medskip
\ifdefined\hypertarget\hypertarget{sourcepointer-2}{}\fi
\noindent\textit{Source reference.} Original reference algebra-lemma-ext-noetherian: sources/stacks-dualizing/pinned-original-source.tex, source lines 1958--1958. Navigation pointer only: the counted text refers to the Stacks item algebra-lemma-ext-noetherian, which belongs to another chapter of the project and is not part of this card. The pointer carries the target's own label as its number because this chapter's staged snapshot does not contain the other chapter's tag table; its content is not reproduced here.\par\medskip
\ifdefined\hypertarget\hypertarget{sourcepointer-3}{}\fi
\noindent\textit{Source reference.} Original reference derived-section-triangulated-sub: sources/stacks-dualizing/pinned-original-source.tex, source lines 1363--1363. Navigation pointer only: the counted text refers to the Stacks item derived-section-triangulated-sub, which belongs to another chapter of the project and is not part of this card. The pointer carries the target's own label as its number because this chapter's staged snapshot does not contain the other chapter's tag table; its content is not reproduced here.\par\medskip
\begin{lemma}
\label{lemma-exact-support-coherent}
With notation as above, assume $A \to B$ is a finite ring map of
Noetherian rings. Then $R\Hom(B, -)$ maps
$D^+_{\textit{Coh}}(A)$ into $D^+_{\textit{Coh}}(B)$.
\end{lemma}

\noindent
Let $A$ be a Noetherian ring. We will denote
$$
D_{\textit{Coh}}(A) \subset D(A)
$$
the full subcategory consisting of those objects $K$ of $D(A)$
whose cohomology modules are all finite $A$-modules. This makes sense
by Derived Categories, Section \ref{derived-section-triangulated-sub}
because as $A$ is Noetherian, the subcategory of finite $A$-modules
is a Serre subcategory of $\text{Mod}_A$.

\begin{proof}
We have to show: if $K \in D^+(A)$ has finite cohomology modules, then the
complex $R\Hom(B, K)$ has finite cohomology modules too.
This follows for example from Lemma \ref{lemma-RHom-ext}
if we can show the ext modules $\Ext^i_A(B, K)$
are finite $A$-modules. Since $K$ is bounded below there is a
convergent spectral sequence
$$
\Ext^p_A(B, H^q(K)) \Rightarrow \text{Ext}^{p + q}_A(B, K)
$$
This finishes the proof as the modules $\Ext^p_A(B, H^q(K))$
are finite by
Algebra, Lemma \ref{algebra-lemma-ext-noetherian}.
\end{proof}

\end{document}
